% Conference-style draft (article class for portability; swap to venue template for submission).
% Build: pdflatex main.tex && bibtex main && pdflatex main.tex && pdflatex main.tex
\documentclass[11pt]{article}
\usepackage[margin=1in]{geometry}
\usepackage{amsmath,amssymb}
\usepackage{booktabs}
\usepackage{microtype}
\usepackage[numbers,sort&compress]{natbib}
\usepackage{hyperref}
\usepackage{url}

\title{Multi-Head Local Spectral Anchor Pooling:\\
Local Time--Frequency Readouts for Protein Language Model Regression}
\author{Anonymous submission}
\date{\today}

\newcommand{\method}{MLSAP}
\newcommand{\R}{\mathbb{R}}

\begin{document}
\maketitle

\begin{abstract}
We study \emph{sequence readout} for antimicrobial-peptide (AMP) activity-cliff regression with frozen protein language models (PLMs).
Standard pooling (mean, max, attention) and recent learnable readouts (multi-layer pooling, latent attention, optimal-transport aggregation) summarize token embeddings without an explicit \emph{local spectral} inductive bias.
We propose \textbf{Multi-Head Local Spectral Anchor Pooling (\method)}, a drop-in operator on last-layer hidden states: per-head analysis projections are transformed with a short-time Fourier transform (STFT); each time--frequency frame is described by magnitude (optionally phase) features; learnable anchors assign frame importance; frame weights are mapped to token weights via overlap-aware accumulation; head outputs are concatenated, optionally gated, and projected.
Across four dataset--backbone settings in our merged benchmark sheet, the best \method{} configuration reaches test Spearman $\rho=0.559$ on \textit{E.\,coli} with ESM2-t12 (below last-token pooling at $\rho=0.563$ on that split) and $\rho=\mathbf{0.462}$ on \textit{S.\,aureus} with ESM2-t12, outperforming sliced-Wasserstein ($0.451$), last-token ($0.443$), and latent-attention ($0.440$) readouts on the same sheet.
We clarify that \method{} summarizes \emph{single-sequence} local structure and does not replace evolutionary multi-sequence information.
\end{abstract}

\section{Introduction}
PLMs such as ESM-2 produce per-token representations of variable length $T$ and hidden size $d$ \cite{lin2023evolutionary}.
Downstream tasks that require one vector per sequence need a \emph{readout} map $\R^{T\times d}\rightarrow \R^{d}$.
Mean pooling remains a strong default; embedding-model surveys systematically compare mean, last-token, and trainable pooling designs \cite{tang2024poolingattentioneffectivedesigns}.
Complementary to token-mixing attention, optimal-transport-style aggregation treats tokens as a set or distribution and can outperform mean pooling on protein tasks \cite{naderi2024aggregating}.

For AMP sequences, functional cues often manifest as short-range physicochemical motifs rather than a single global statistic along the chain.
We therefore introduce \emph{local} time--frequency structure before aggregation: STFT frames yield descriptors that are scored against learnable anchors, producing token weights without interpolating a \emph{global} frequency axis back to residues.

\paragraph{Contributions.}
\begin{itemize}
  \item We formulate \method{}, a multi-head STFT--anchor readout with overlap-aware frame-to-token mapping, implemented as \texttt{local\_spectral\_anchor} in our codebase.
  \item We position \method{} as a \emph{drop-in} replacement for sequence pooling on PLM hidden states, alongside mean/max/attention, latent attention, multi-layer trainable pooling (MLTP), and OT-style sliced-Wasserstein readouts.
  \item We report test Spearman correlations from a merged spreadsheet of grid and toggle ablations on two AMP datasets and two ESM2 backbones (t6/t12), with fixed DC-processing branch (\texttt{apply=none}) for fair readout comparison where applicable.
\end{itemize}

\section{Related work}
\paragraph{Pooling for PLM embeddings.}
\citet{tang2024poolingattentioneffectivedesigns} survey pooling and attention choices for LLM-based embedding models, including mean and last-token readouts and trainable pooling layers.
Our experiments include analogous readouts on top of ESM2.

\paragraph{Distribution-based aggregation.}
\citet{naderi2024aggregating} aggregate residue-level PLM embeddings with optimal transport and sliced-Wasserstein distances, improving over mean pooling on several protein tasks.
Our implementation uses an OT-inspired sliced-Wasserstein readout (\texttt{swe\_ot}) as a strong baseline.

\paragraph{Spectral views of sequences.}
Time--frequency transforms expose local periodicities compared with a single global Fourier spectrum.
We do \emph{not} claim cross-modal universality; we evaluate \method{} on AMP PLM hidden states only.
Related literature on wavelet or STFT attention in other domains is omitted pending verified bibliography entries.

\section{Method}
Let $X\in\R^{B\times T\times D}$ denote last hidden states from the PLM, with attention mask $M\in\{0,1\}^{B\times T}$.
We split channels into $H$ heads with $D=H\cdot d_h$ and process each head independently.

\paragraph{Analysis projection and STFT.}
For head $h$, $X_h\in\R^{B\times T\times d_h}$ is mapped to $A_h\in\R^{B\times T\times p}$ with a linear map (analysis dimension $p$).
Channels are reshaped and fed to \texttt{torch.stft} with Hann window, \texttt{n\_fft}, \texttt{win\_length}, \texttt{hop\_length}, optional centering, and onesided magnitude (or real/imaginary concatenation if phase is enabled).
The result is reshaped to frame descriptors $F_h\in\R^{B\times F\times d_f}$, where $F$ is the number of STFT frames and $d_f$ depends on $p$, frequency bins, and phase mode.

\paragraph{Anchor assignment.}
Learnable anchors $A^{(h)}\in\R^{K\times d_f}$ per head induce soft assignments by scaled Euclidean distance to projected frames, followed by softmax with a per-head temperature.
We average anchor probabilities per frame to obtain a scalar frame weight vector.

\paragraph{Frame-to-token mapping.}
Each frame $m$ contributes its scalar weight uniformly over token indices covered by the STFT support (with an optional left padding heuristic when centering is enabled); contributions are divided by per-token coverage, multiplied by the mask, and normalized to sum to one per sequence; degenerate rows fall back to uniform or masked-uniform weights.

\paragraph{Aggregation and head fusion.}
Token weights yield a weighted sum of $X_h$ over time; concatenating heads gives $\R^{BD}$; optional sigmoid gating and a linear output projection return the pooled vector in $\R^{D}$.

\paragraph{Scope and limitations.}
Anchors live in \emph{frame descriptor} space, not in a global frequency bin axis mapped to residues.
The mapping is not an inverse STFT; it is overlap-aware averaging, not Hann-weighted overlap-add reconstruction.
When \texttt{center=True}, padding alignment may differ slightly from \texttt{torch.stft} across PyTorch versions.
\method{} uses only the input sequence; it does not encode evolutionary information from multiple homologs.

\section{Experiments}
\paragraph{Task and metrics.}
We report test Spearman $\rho$ (primary) and Pearson $r$ from training runs logged in Hydra output directories and summarized in the merged workbook \texttt{local\_spectral\_anchor\_grid\_2x2\_with\_pooling\_mean\_max\_attn\_merged.xlsx} under \texttt{outputs/2026-04-03/}.
Metrics are computed on held-out test splits defined by the project configuration (activity-cliff regression on \textit{E.\,coli} and \textit{S.\,aureus} datasets).

\paragraph{Baselines.}
\textbf{Classic pooling} (\texttt{pooling\_baseline}): mean, max, and attention where present in the merged sheet.
\textbf{SpectralAnchorV2 internal readouts} (\texttt{spectral\_anchor\_v2\_ablation}): mean/max/attention variants bundled in the same ablation sheet---distinct from \method{}.
\textbf{LLM readout baselines} (\texttt{llm\_readout\_pooling}): last-token, latent attention, \texttt{swe\_ot}, and \texttt{mltp}, all with \texttt{apply=none} in the logged run names.

\paragraph{MLSAP grid and toggles.}
We sweep head/anchor counts, analysis width, \texttt{n\_fft}, hop length, and toggles for STFT centering, phase, and gating (see \texttt{config\_name} strings prefixed with \texttt{local\_spectral\_anchor/} in the workbook).
The \textbf{best MLSAP} row per sheet maximizes test Spearman among \method{} ablations.

% Table generated from outputs/2026-04-03/local_spectral_anchor_grid_2x2_with_pooling_mean_max_attn_merged.xlsx
% Metrics: test Spearman correlation (higher is better). See paper/main.tex for discussion.
\begin{table*}[t]
  \centering
  \small
  \caption{Test Spearman correlation ($\rho$) on AMP activity-cliff regression. All readouts use the same frozen PLM backbone with \texttt{apply=none} unless noted. Best MLSAP denotes the strongest \texttt{local\_spectral\_anchor} row from the local $2\times2$ grid (heads $\times$ anchors) and toggle ablations in the merged workbook. ``---'' indicates missing baseline rows in the merged sheet for that backbone/dataset.}
  \label{tab:main}
  \begin{tabular}{lcccc}
    \toprule
    Readout / pooling & \multicolumn{2}{c}{ESM2-t12} & \multicolumn{2}{c}{ESM2-t6} \\
    \cmidrule(lr){2-3} \cmidrule(lr){4-5}
    & \textit{E.\,coli} & \textit{S.\,aureus} & \textit{E.\,coli} & \textit{S.\,aureus} \\
    \midrule
    Mean / max / attn (classic) & 0.531 / 0.537 / --- & --- / --- / --- & 0.532 / 0.536 / 0.523 & 0.394 / 0.414 / 0.438 \\
    SpectralAnchorV2 (mean / max / attn) & 0.504 / 0.553 / 0.555 & 0.388 / 0.427 / 0.406 & --- & --- \\
    Last-token & \textbf{0.563} & 0.443 & 0.517 & 0.389 \\
    Latent attention & 0.535 & 0.440 & 0.534 & 0.424 \\
    Sliced-Wasserstein / OT-style & 0.558 & 0.451 & 0.536 & 0.436 \\
    Multi-layer trainable pooling & 0.541 & 0.418 & 0.520 & 0.334 \\
    Best MLSAP (ours) & 0.559 & \textbf{0.462} & \textbf{0.540} & \textbf{0.464} \\
    \bottomrule
  \end{tabular}
\end{table*}


\paragraph{Observations.}
On \textit{E.\,coli} with ESM2-t12, last-token pooling achieves the highest $\rho$ among reported readouts ($0.563$), slightly above the best \method{} ($0.559$).
On \textit{S.\,aureus} with ESM2-t12, \method{} ($0.462$) outperforms OT-style ($0.451$), last-token ($0.443$), and latent-attention ($0.440$) baselines on the same sheet.
ESM2-t6 shows best \method{} $\rho=0.540$ on \textit{E.\,coli}, ahead of OT-style ($0.536$) and mean pooling ($0.532$); on \textit{S.\,aureus}, \method{} ($0.464$) exceeds attention pooling ($0.438$) and OT-style ($0.436$), with MLTP lower ($0.334$) in this merge.

\section{Limitations}
\method{} depends on STFT hyperparameters and may be sensitive to sequence length and padding; short sequences are padded for numerical stability in STFT, while token mapping still uses the original length $T$.
We compare single runs per configuration in the merged sheet; confidence intervals and multi-seed variance are not reported here.
Some backbone--dataset combinations lack full classic pooling rows in the merged workbook (e.g., ESM2-t12 classic baselines for \textit{S.\,aureus}); results should be interpreted conditional on available rows.

\section{Conclusion}
We introduced \method{}, a multi-head local STFT--anchor pooling operator for PLM token embeddings, and evaluated it against strong readout baselines on AMP regression.
The best configuration is competitive with or better than several learnable baselines on multiple splits, while last-token remains strongest on one ESM2-t12 \textit{E.\,coli} setting in our table.
Future work includes statistical aggregation across seeds, tighter alignment between STFT centering and token support, and verified citations for spectral sequence models.

\bibliographystyle{plainnat}
\bibliography{references}

\end{document}
